\documentclass[10pt]{mypackage}

\usepackage[uprightgreeks,light]{kpfonts}
\renewcommand{\coloneq}{\coloneqq}
\usepackage{homework}
%\usepackage{notes}

%\usepackage[ backend=bibtex, style = alphabetic, sorting=ynt ]{biblatex}
%\addbibresource{  }

\usepackage{parskip}

\fancyhf{}
\fancyhead[R]{Avinash Iyer}
\fancyhead[L]{Algebra II: Homework 8}
\fancyfoot[C]{\thepage}

\setcounter{secnumdepth}{0}

\begin{document}
\RaggedRight
\begin{problem}[Problem 1]
  Given a ring $R$ with $1$, the opposite ring $R^{\op}$ is defined as follows. As a group with addition, $R^{\op}$ is the same as $R$, and multiplication $\ast$ in $R^{\op}$ is defined by $a\ast b = ba$.
  \begin{enumerate}[(a)]
    \item Let $\operatorname{end}_R(R)$ be the ring of endomorphisms of $R$ considered as a left module over itself. Prove that $\operatorname{end}_R(R)$ is isomorphic to $R^{\op}$.
    \item Let $G$ be a group, $K$ a commutative ring with $1$, and $K[G]$ the corresponding group ring. Prove that $K[G]$ is isomorphic to its opposite ring.
  \end{enumerate}
\end{problem}
\begin{solution}\hfill
  \begin{enumerate}[(a)]
    \item Let $a\in R$, and let $\varphi\in \operatorname{end}_R(R)$. Since $R$ has $1$ and $\varphi$ is a left $R$-module endomorphism, it follows that
      \begin{align*}
        \varphi\left( a \right) &= a\varphi\left( 1 \right),
      \end{align*}
      meaning that $\varphi$ is determined by where $1$ is mapped to. Call this value $b$, label $\varphi = \varphi_b$, and consider the map $b\xmapsto{\Phi} \varphi_b$ mapping from $R^{\op}$ to $\operatorname{end}_R(R)$. We claim this is an isomorphism of rings. First, to show it is a homomorphism, observe that if $b,c\in R$, then
      \begin{align*}
        \varphi_b\circ \varphi_c(1) &= \varphi_b\left(c\right)\\
                                    &= c\varphi_b(1)\\
                                    &= cb\\
                                    &= b\ast c,
      \end{align*}
      where the operation $\ast$ is the opposite operation in $R^{\op}$. In particular, since the map is also additive by construction, it follows that this is a homomorphism.

      Now, note that $\varphi_b$ is the zero map if and only if $ab = 0$ for all $a\in R$, which occurs only when $b = 0$. Therefore, the map $\Phi$ is injective. Surjectivity follows from the fact that we may define a right inverse $\Psi\colon \operatorname{end}_R(R)\rightarrow R^{\op}$ taking $\varphi\mapsto \varphi(1)$, satisfying $\Phi(\Psi(\varphi)) = \Phi(\varphi(1)) = \varphi_{\varphi(1)} = \varphi$.
  \end{enumerate}
\end{solution}
\begin{problem}[Problem 2]
  Let $p$ be a prime, and let $G = \operatorname{Heis}\left( \F_p \right)$ be the Heisenberg group over $\F_p$.
  \begin{enumerate}[(a)]
    \item Determine the number of conjugacy classes of $G$ and their sizes.
    \item Let $\omega\neq 1$ be a $p$th root of unity, that is $\omega = e^{2\pi i k/p}$ with $1\leq p \leq p=1$. Let $V$ be a $p$-dimensional complex vector space with basis $e_{[0]},\dots,e_{[p-1]}$, where we think of indices as elements of $\F_p$. Prove that there exists a representation $\left( \rho_{\omega},V \right)$ for which
      \begin{itemize}
        \item $\rho_{\omega}(z)e_{[k]} = \omega e_{[k]}$;
        \item $\rho_{\omega}(y)e_{[k]} = e_{[k+1]}$;
        \item $\rho_{\omega}(x)e_{[k]} = \omega^{k}e_{[k]}$.
      \end{itemize}
  \end{enumerate}
\end{problem}
\begin{solution}\hfill
  \begin{enumerate}[(a)]
    \item Fix the presentation
      \begin{align*}
        G &= \left\langle x,y,z|\set{x^{p},y^{p},z^{p},xzx^{-1}z^{-1},yzy^{-1}z^{-1},xyx^{-1}y^{-1}z^{-1}} \right\rangle.
      \end{align*}
      Observe that $Z(G) = \left\langle z \right\rangle$, since we know that $G$ is a group of order $p^{3}$ that is not abelian (hence its center has order $p$), and $z$ is central. In particular, this means that in the class formula
      \begin{align*}
        \left\vert G \right\vert &= \left\vert Z(G) \right\vert + \sum_{i=1}^{r} \left[ G:Z_G\left(g_i\right) \right]\\
        p^3 &= p + \sum_{i=1}^{r} \left[ G:Z_G\left( g_i \right) \right],
      \end{align*}
      we have that the sum of the sizes of the conjugacy classes with representatives $g_1,\dots,g_r$ is given by $p^3 - p$. Yet, since each $\left[ G:Z_G\left( g_i \right) \right]$ must divide $p^3$, we have that the factoring $p^3 - p = p\left( p^2 - 1 \right)$ implies that there are $p^2 - 1$ conjugacy classes each with size $p$.
    \item 
  \end{enumerate}
\end{solution}
\begin{problem}[Problem 3]
  Let $G = S_n$ for some $n\geq 2$, and let $\chi$ be an irreducible complex character of $G$. Prove that $\chi$ is real-valued; that is, $\chi(g) \in \R$ for all $g\in G$.
\end{problem}
\begin{solution}
  Let $\rho$ be the representation whose character is $\chi$. If $\rho\colon G\rightarrow V$ is the given finite-dimensional irreducible (unitary) representation, then we observe that it is sufficient to show that
  \begin{align*}
    \chi_{\rho^{\ast}}\left( g \right) &= \chi\left( g \right)
  \end{align*}
  for all $g\in G$, as we know that $\chi_{\rho^{\ast}}\left( g \right) = \overline{\chi\left( g \right)}$. Since $g\in G$ is some permutation, $g$ decomposes into a product of disjoint cycles given by
  \begin{align*}
    g &= \sigma_1\sigma_2\cdots \sigma_k
  \end{align*}
  for some $k\leq n$. Since these cycles are disjoint, they commute, and since the inverse of a cycle with a given length is a cycle of the same length, we observe then that
  \begin{align*}
    g^{-1} &= \sigma_1^{-1}\sigma_2^{-1}\cdots \sigma_k^{-1}.
  \end{align*}
  In particular, both $g$ and $g^{-1}$ have the same cycle type, meaning that they have the same character. Yet, this gives
  \begin{align*}
    \chi_{\rho^{\ast}}\left( g \right) &= \Tr \left( \left[ \rho^{\ast}\left( g \right) \right]_{\beta^{\ast}} \right)\\
                                       &= \Tr \left( \left[ \rho\left( g^{-1} \right) \right]_{\beta}^{T} \right)\\
                                       &= \Tr \left( \left[ \rho\left( g^{-1} \right) \right]_{\beta} \right)\\
                                       &= \Tr \left( \left[ \rho\left( g \right) \right]_{\beta} \right)\\
                                       &= \chi\left( g \right).
  \end{align*}
  Therefore, $\chi(g)\in \R$.
\end{solution}
\begin{problem}[Problem 4]
  Give an example of two representations $V$ and $W$ of the same group which are not equivalent but have the same dimension and the same character.
\end{problem}
\begin{solution}
  Let $G = \Z/2\Z$, and consider the representation $\rho\colon G\rightarrow \GL_2\left(\F_2\right)$ given by
  \begin{align*}
    0 &\mapsto \begin{pmatrix}1 & 0 \\ 0 & 1\end{pmatrix}\\
    1 &\mapsto \begin{pmatrix}0 & 1 \\ 1 & 0\end{pmatrix}.
  \end{align*}
  This is a representation since
  \begin{align*}
    \begin{pmatrix}0 & 1 \\ 1 & 0\end{pmatrix} \begin{pmatrix}0 & 1 \\ 1 & 0\end{pmatrix} &= \begin{pmatrix}1 & 0 \\ 0 & 1\end{pmatrix},
  \end{align*}
  and it has dimension $2$ and character $0$ for both $\rho(0)$ and $\rho(1)$. Meanwhile, if $\1\colon G\rightarrow \GL_2\left( \F_2 \right)$ is the representation taking both $0$ and $1$ to the identity matrix, it follows that $\1$ and $\rho$ are not equivalent since the identity representation is unique. Yet, $\1$ has character $0$ and dimension $2$.
\end{solution}
\end{document}
